\section{Justification and Formal Status of the Axioms}
\label{app:axiom_justifications}

This appendix provides the rationale for each of the twelve axioms of the \Apeiron{} Framework and, where applicable, formal proofs within specific mathematical frameworks.

\subsection*{Axiom 1 (Irreducibility)}
\textbf{Rationale:} Without a minimal granularity, any representational system succumbs to infinite regress. This axiom is not provable in an absolute sense but is a methodological necessity. In specific domains, however, it can be instantiated. For example, in the Boolean algebra of natural numbers under divisibility, the integer $1$ is provably an atom (Proposition~\ref{prop:bool-atom}). In the category of pointed sets, the singleton is provably a zero object (Proposition~\ref{prop:cat-atom}).

\subsection*{Axiom 2 (Multi-Axial Atomicity)}
\textbf{Rationale:} No single mathematical framework captures all aspects of irreducibility. Convergence across independent frameworks provides robustness against the limitations of any one formalism. The falsification experiment in Section~\ref{sec:validation} demonstrates that for the integer $1$, all four frameworks converge.

\subsection*{Axiom 3 (Relational Constitution)}
\textbf{Rationale:} This axiom operationalizes structural realism and the insight from category theory that objects are defined by their morphisms. It cannot be proven a priori but is a design choice that distinguishes the framework from substance-based ontologies.

\subsection*{Axiom 4 (Emergence)}
\textbf{Rationale:} The irreducibility of higher-level phenomena to lower-level constituents is a well-documented empirical fact in complex systems \cite{bedau1997, mitchell2009}. The axiom formalizes this as an architectural commitment.

\subsection*{Axiom 5 (Temporal Endogeneity)}
\textbf{Rationale:} In relational quantum mechanics \cite{rovelli1996} and general relativity, time is not a background parameter but emerges from correlations. The axiom aligns the framework with these physical theories and avoids hard-coding a Newtonian time axis.

\subsection*{Axiom 6 (Ontological Openness)}
\textbf{Rationale:} Any fixed set of categories would reintroduce anthropocentric bias. The axiom ensures that the system can extend its own ontology, a capacity demonstrated empirically in large language models \cite{wei2022}.

\subsection*{Axiom 7 (Autopoietic Closure)}
\textbf{Rationale:} Autopoiesis is the defining characteristic of living systems \cite{maturana1980}. The axiom ensures that world-models are self-maintaining, a necessary condition for genuine world-making.

\subsection*{Axiom 8 (Ethical Reversibility)}
\textbf{Rationale:} This is a normative axiom, not a descriptive one. It encodes the precautionary principle and ensures that the framework is aligned with human values. It can be formalized as an $\varepsilon$-reversibility condition (Definition~14.1).

\subsection*{Axiom 9 (Observer Relativity)}
\textbf{Rationale:} All measurements and judgments are relative to an observer's framework \cite{barad2007}. The \texttt{MetaSpecification} makes this relativity explicit and computationally tractable.

\subsection*{Axiom 10 (Generativity)}
\textbf{Rationale:} An observable that is merely atomic but lacks connections to other observables cannot serve as a substrate for emergence. The axiom quantifies this capacity, and its components (resonance, relational degree, consensus) are measurable.

\subsection*{Axiom 11 (Epistemic Humility)}
\textbf{Rationale:} No system can fully certify its own correctness (G\"odel). The axiom mandates that the system maintain explicit representations of uncertainty, enabling human oversight and iterative improvement.

\subsection*{Axiom 12 (Convergence)}
\textbf{Rationale:} The framework is designed to approach a fixed point of self-consistency under iterative refinement. This is a regulative ideal, analogous to the concept of truth as a limit in Peircean semiotics. In specific cases, such as the consensus update in Layer~17, convergence can be proven under contraction conditions (Theorem~A.1).